\documentclass[11pt,a4paper]{article}
\usepackage[T1]{fontenc}
\usepackage[utf8]{inputenc}
\usepackage[french]{babel}
\usepackage{lmodern}
\usepackage[margin=2.2cm]{geometry}
\usepackage{amsmath,graphicx,booktabs,xcolor,caption}
\usepackage{pgfplots}
\pgfplotsset{compat=1.16}
\tikzset{every picture/.prefix code={\shorthandoff{;:!?}}}
\usepackage[hidelinks]{hyperref}
\setlength{\parindent}{0pt}\setlength{\parskip}{6pt}

\definecolor{frozen}{HTML}{1552AD}
\definecolor{full}{HTML}{C13B21}
\newcommand{\src}{../../../../reports/experiments/experiment_81_execution_20261106/figures}
\newcommand{\lecture}[1]{\textbf{Lecture.} #1}

% Les quatre courbes d'une grandeur : couleur = bras, trait = jeu de vidéos.
\newcommand{\quatre}[1]{%
  \addplot[frozen,thick] table[x=update,y=#1]{frozen_train.dat};
  \addplot[frozen,thick,densely dashed] table[x=update,y=#1]{frozen_selection.dat};
  \addplot[full,thick] table[x=update,y=#1]{full_train.dat};
  \addplot[full,thick,densely dashed] table[x=update,y=#1]{full_selection.dat};}
\pgfplotsset{courbe/.style={width=\linewidth,height=5.3cm,xmin=0,xmax=100,
  xlabel={Mise à jour de l'optimiseur},grid=major,grid style={gray!20},
  extra x ticks={20.5},extra x tick labels={},extra x tick style={grid=major,grid style={black!60,densely dotted,thick}},
  tick label style={font=\footnotesize},label style={font=\footnotesize},title style={font=\small\bfseries}}}

\begin{document}

{\LARGE\bfseries Expérience 81 : lecture des figures}\\[2pt]
{\small Document du 10 novembre 2026. Aucun entraînement n'a été relancé pour ce document.}

\section*{Ce qui a été entraîné}

80 vidéos thermiques : 40 sans déversement (\emph{clean}) et 40 avec déversement (\emph{cold}).
64 servent à l'entraînement et 16 à la sélection. Deux versions du même modèle sont comparées,
avec les mêmes vidéos, la même graine (239) et 100 mises à jour chacune :

\begin{itemize}\itemsep1pt
\item \textcolor{frozen}{\textbf{frozen} (bleu)} : l'encodeur vidéo pré-entraîné V-JEPA~2 reste figé ; seul le petit module temporel est entraîné.
\item \textcolor{full}{\textbf{full} (rouge)} : les 24 blocs de l'encodeur sont aussi entraînés.
\end{itemize}

La fonction de coût est la somme de quatre termes :
\[
L \;=\; L_{\text{ViCo}} \;+\; 0{,}1\,L_{\text{align}} \;+\; L_{\text{predict}} \;+\; 0{,}1\,L_{\text{keep}} .
\]

Dans toutes les courbes : trait plein = vidéos d'entraînement, trait tireté = vidéos de sélection.
La ligne verticale pointillée marque la fin du curriculum de 20 mises à jour ; l'état de
l'optimiseur AdamW est remis à zéro avant la mise à jour 21, comme le prévoit le protocole.

\textbf{Sur les courbes redessinées.} Les figures 2 et 3 sont retracées à partir des valeurs
enregistrées (\texttt{figures/loss\_history.csv}), sans lissage ni modification d'aucune valeur.
Seule la présentation change : axes gradués, légende, échelle logarithmique quand la courbe
couvre plusieurs ordres de grandeur. Les figures d'origine sont reproduites en annexe.

\textbf{Limite à garder en tête.} Les 16 vidéos de sélection viennent des mêmes conditions
d'acquisition que les vidéos d'entraînement et ont servi à de nombreuses vérifications. Ce n'est
pas un jeu de test indépendant.

\newpage
\section*{Figure 1 : t-SNE avant tout entraînement}

\begin{center}\includegraphics[width=.8\linewidth]{\src/tsne_frozen_initial.png}\end{center}

\textbf{Ce qui est tracé.} Chaque point est une vidéo, représentée par le vecteur de 128
dimensions que produit le modèle à la mise à jour 0, c'est-à-dire avant tout apprentissage.
Le t-SNE projette ces 80 vecteurs en deux dimensions en gardant proches les vidéos semblables.
Bleu = sans déversement, rouge = avec déversement ; rond = entraînement (64), carré = sélection (16).

\textbf{Ce qu'on observe.}
\begin{itemize}\itemsep1pt
\item Les points bleus et rouges forment deux groupes disjoints ; aucun point d'une classe
      ne se trouve dans le groupe de l'autre.
\item Un petit groupe bleu isolé à gauche contient 7 vidéos : \texttt{videoFrame} 1, 3, 4, 5, 6, 7 et 8
      (\texttt{figures/tsne\_frozen\_initial.json}).
\item À cette même mise à jour 0, les 16 vidéos de sélection sont déjà toutes classées
      correctement, et 63 des 64 vidéos d'entraînement (\texttt{frozen/history.json}).
\end{itemize}

\lecture{Les deux classes sont séparées avant que le modèle ait appris quoi que ce soit.
La séparation ne vient donc pas de la fonction de coût. Elle est déjà dans les vidéos elles-mêmes :
les vidéos des deux classes ont été filmées dans des conditions différentes (384$\times$288 pour
les unes, 640$\times$480 pour les autres, scènes différentes). Cette explication est une inférence :
la figure montre que la séparation préexiste, pas ce qui la cause.}

\textbf{Ce que la figure ne montre pas.} Les axes et les distances entre groupes d'un t-SNE n'ont pas
de signification, et deux t-SNE ne se comparent pas entre eux. La figure ne dit rien de l'endroit
que le modèle regarde dans l'image.

\newpage
\section*{Figure 2 : les termes de la fonction de coût}

\begin{center}
\begin{tikzpicture}
\begin{axis}[hide axis,xmin=0,xmax=1,ymin=0,ymax=1,width=3cm,height=2cm,
  legend columns=4,legend style={draw=none,font=\footnotesize,/tikz/every even column/.append style={column sep=10pt}}]
\addlegendimage{frozen,thick}\addlegendentry{frozen, entraînement}
\addlegendimage{frozen,thick,densely dashed}\addlegendentry{frozen, sélection}
\addlegendimage{full,thick}\addlegendentry{full, entraînement}
\addlegendimage{full,thick,densely dashed}\addlegendentry{full, sélection}
\end{axis}
\end{tikzpicture}

\begin{minipage}{.49\linewidth}
\begin{tikzpicture}\begin{semilogyaxis}[courbe,title={(a) ViCo, terme de classe (échelle log)},ymin=1e-8,ymax=5]
\quatre{semantic}\end{semilogyaxis}\end{tikzpicture}
\end{minipage}\hfill
\begin{minipage}{.49\linewidth}
\begin{tikzpicture}\begin{axis}[courbe,title={(b) Prédiction du futur},ymin=0,ymax=0.35,
  yticklabel style={/pgf/number format/fixed,/pgf/number format/precision=2}]
\quatre{prediction}
\addplot[black,thin] table[x=update,y=persistence]{full_selection.dat};
\node[font=\scriptsize,anchor=south east] at (axis cs:99,0.039) {persistance (sélection)};
\end{axis}\end{tikzpicture}
\end{minipage}

\begin{minipage}{.49\linewidth}
\begin{tikzpicture}\begin{axis}[courbe,title={(c) Rétention},ymin=0,ymax=0.35,
  yticklabel style={/pgf/number format/fixed,/pgf/number format/precision=2}]
\quatre{retention}\end{axis}\end{tikzpicture}
\end{minipage}\hfill
\begin{minipage}{.49\linewidth}
\begin{tikzpicture}\begin{axis}[courbe,title={(d) Alignement},ymin=0,ymax=0.003,
  scaled y ticks=base 10:3]
\quatre{alignment}\end{axis}\end{tikzpicture}
\end{minipage}
\end{center}
{\footnotesize Valeurs brutes, avant application des poids 0,1. Source : \texttt{figures/loss\_history.csv}.
Évaluation à chaque mise à jour avec les poids finaux du curriculum, pour que les points soient comparables.}

\textbf{(a) ViCo, le terme de classe.} Il est faible quand les vidéos de même classe ont des
représentations proches et celles de classes opposées des représentations éloignées. C'est le seul
terme qui utilise les étiquettes.
\begin{itemize}\itemsep1pt
\item \textcolor{full}{full} : il tombe sous $10^{-3}$ après la mise à jour 12 (sélection) et 18 (entraînement).
\item \textcolor{frozen}{frozen} : il oscille pendant le curriculum, remonte à 2,27 (entraînement) et 2,45
      (sélection) à la mise à jour 21, juste après la remise à zéro de l'optimiseur, puis repasse sous
      $10^{-3}$ après la mise à jour 32. La cause exacte de ce pic n'a pas été analysée.
\item Valeurs finales sur la sélection : $2{,}3\times10^{-7}$ (frozen) et $2{,}2\times10^{-8}$ (full).
\end{itemize}
\lecture{Le problème de classification est résolu en quelques dizaines de mises à jour. Une fois ce
terme proche de zéro, il ne fournit presque plus de signal d'apprentissage : rien, dans ce terme, ne
pousse ensuite le modèle à chercher le liquide.}

\textbf{(b) Prédiction du futur.} Erreur du modèle quand il prédit les caractéristiques de la vidéo
1, 2 et 4 pas de temps plus tard. La ligne noire est la \emph{persistance} : l'erreur obtenue en
supposant simplement que rien ne change (0,0391 sur la sélection). Ce n'est pas un terme entraîné,
c'est le repère à battre.
\begin{itemize}\itemsep1pt
\item L'erreur monte d'abord jusqu'à environ 0,33 (sélection, mise à jour 2), puis décroît.
\item Elle passe sous la persistance à partir de la mise à jour 28 (full) et 49 (frozen), sur la sélection.
\item Valeurs finales sur la sélection : 0,0244 (frozen, 38\,\% sous la persistance) et 0,0142
      (full, 64\,\% sous la persistance).
\end{itemize}
\lecture{Le modèle prédit mieux que « rien ne change », surtout quand l'encodeur est entraîné.
La pente est encore légèrement négative sur les 20 dernières mises à jour
(\texttt{measured\_results.json}) : l'apprentissage n'est pas complètement stabilisé. Battre la
persistance ne prouve pas que le modèle suit le mouvement du liquide.}

\textbf{(c) Rétention.} Écart entre les caractéristiques adaptées par le modèle et celles de
l'encodeur pré-entraîné d'origine. Ce terme empêche le modèle d'oublier ce que l'encodeur savait.
\begin{itemize}\itemsep1pt
\item Il vaut 0 au départ, monte jusqu'à 0,31 à 0,33 (sélection, mises à jour 2 et 3), puis redescend.
\item Valeurs finales : 0,0165 / 0,0309 (frozen, entraînement / sélection) et 0,0147 / 0,0154 (full).
\end{itemize}
\lecture{Le modèle s'écarte d'abord des caractéristiques d'origine pour séparer les classes, puis
y revient. Dans le bras frozen, l'écart reste environ deux fois plus grand sur la sélection que sur
l'entraînement.}

\textbf{(d) Alignement.} Écart entre les états temporels de deux vues augmentées de la même vidéo.
L'échelle est environ cent fois plus petite que celle des panneaux (b) et (c) (maximum 0,0028), et le terme est
encore multiplié par 0,1 dans la somme.
\begin{itemize}\itemsep1pt
\item \textcolor{frozen}{frozen} : bosse jusqu'à 0,0028 à la mise à jour 49, puis retour vers 0,0011 à 0,0013.
\item \textcolor{full}{full} : décroissance régulière jusqu'à environ 0,0002.
\end{itemize}
\lecture{Les deux vues d'une même vidéo restent cohérentes. Ce terme pèse moins de 0,0003 dans le
total ; ses variations ne se voient pas sur la courbe totale.}

\newpage
\section*{Figure 3 : la fonction de coût totale}

\begin{center}
\begin{tikzpicture}
\begin{semilogyaxis}[courbe,width=.95\linewidth,height=6.3cm,ymin=0.008,ymax=4,
  ylabel={$L$ (échelle log)},legend style={font=\footnotesize,draw=gray!40},legend pos=north east,
  log ticks with fixed point]
\quatre{total}
\legend{{frozen, entraînement},{frozen, sélection},{full, entraînement},{full, sélection}}
\end{semilogyaxis}
\end{tikzpicture}
\end{center}
{\footnotesize $L = L_{\text{ViCo}} + 0{,}1\,L_{\text{align}} + L_{\text{predict}} + 0{,}1\,L_{\text{keep}}$.
Source : \texttt{figures/loss\_history.csv}.}

\textbf{Ce qui est tracé.} La somme pondérée des quatre termes de la figure 2, pour les deux bras et
les deux jeux de vidéos.

\textbf{Ce qu'on observe.}
\begin{itemize}\itemsep1pt
\item Au début, le total suit le terme ViCo, qui domine la somme : mêmes oscillations pendant le
      curriculum, et même pic du bras frozen à la mise à jour 21 (2,35 en entraînement, 2,56 en sélection).
\item Une fois ViCo proche de zéro, le total vaut pratiquement
      $L_{\text{predict}} + 0{,}1\,L_{\text{keep}}$ et décroît lentement.
\item Le bras full descend plus vite et plus bas.
\end{itemize}

\begin{center}\small
\begin{tabular}{lcc}\toprule
Valeur finale du total (mise à jour 100) & Entraînement & Sélection\\\midrule
frozen & 0,0260 & 0,0276\\
full   & 0,0124 & 0,0157\\\bottomrule
\end{tabular}
\end{center}

\lecture{L'optimisation fonctionne dans les deux bras, et les courbes d'entraînement et de sélection
restent proches. Cela décrit la qualité de l'optimisation, pas ce que le modèle a appris à regarder.}

\section*{Ce que ces trois figures établissent, et ce qu'elles n'établissent pas}

\begin{tabular}{p{.47\linewidth}p{.47\linewidth}}\toprule
\textbf{Établi (mesuré)} & \textbf{Non établi}\\\midrule
Les deux classes sont séparées avant tout entraînement (figure 1). &
Que cette séparation repose sur le liquide.\\[3pt]
Le terme de classe tombe à presque zéro en moins de 35 mises à jour (figure 2a). &
Que le modèle porte son attention sur le déversement.\\[3pt]
Les 16 vidéos de sélection sont classées correctement avant et après l'entraînement, dans les deux bras. &
Que le modèle fonctionnerait dans une scène nouvelle.\\\bottomrule
\end{tabular}

Aucune de ces courbes ne mesure l'endroit où le modèle regarde. Dans cette même expérience, les
cartes d'attention ont été examinées séparément : le modèle classe correctement, mais ne se focalise
pas de façon fiable sur le liquide dans les trois vidéos connues pour poser problème (cold56, cold10
et cold13 ; \texttt{report.md}). Les deux constats sont compatibles : puisque la scène suffit à
reconnaître la classe, la fonction de coût peut atteindre zéro sans que le liquide soit utilisé.

\newpage
\section*{Annexe : figures d'origine}

Figures produites par l'expérience le 6 novembre 2026, non modifiées. Les courbes des figures 2 et 3
ci-dessus reprennent exactement ces valeurs.

\begin{center}
\includegraphics[width=.9\linewidth]{\src/loss_components.png}\\[8pt]
\includegraphics[width=.8\linewidth]{\src/total_loss.png}
\end{center}

{\footnotesize Dossier source :
\texttt{reports/experiments/experiment\_81\_execution\_20261106/}.}

\end{document}
